\chapter*{Abstract}
%\addcontentsline{toc}{chapter}{Abstract}
\label{chap:abstract}
An intelligent tutor can accurately determine what a student has mastered yet still lose their engagement, since a major driver of disengagement and dropout is affective rather than cognitive, and few tutoring policies treat affect as a control signal. This thesis addresses this gap by developing a two-stage affect-aware tutoring architecture whose perception and decision components are validated independently. The perception module performs multi-label academic-affect recognition from ten-second webcam clips, predicting boredom, engagement, confusion, and frustration as four binary states that may co-occur within the same clip. It combines an ImageNet-pretrained ResNet-50 backbone, Squeeze-and-Excitation channel recalibration, a two-layer bidirectional LSTM, and temporal attention pooling over 24 frames into four sigmoid heads. The decision module formulates tutoring as a partially observable Markov decision process solved by a Deep Q-Network (DQN), with state incorporates affective channels and latent mastery, while a multi-objective reward function jointly optimises Hake’s normalised learning gain and affect regulation across eight pedagogical actions, while the environment dynamics are derived from established theoretical models. To address of incomparable evaluation in prior DAiSEE-based work, six competing architectures were retrained from scratch under a unified experimental pipeline, with class imbalance handled consistently at the data, loss, and decision-threshold levels. The policy was trained on 10,000 heterogeneous synthetic learners and evaluated over ten random seeds against both a rule-based expert tutor and a random controller. The proposed model achieves the highest Macro F1 among standardised configurations (0.453), with a Weighted F1 of 0.798 and a mean per-label binary accuracy of 0.782 across the four affect labels. In simulation, the DQN policy outperforms both baselines across all metrics in both populations, improving learning gain by 0.088 over the expert tutor on the mixed population ($d = 2.01$, $p_{\mathrm{Holm}} < 10^{-2}$, effect sizes ranging from 1.81 to 3.30 overall). It also shows interpretable behavioural patterns and demonstrates that a policy trained on a mixed population generalises effectively to the canonical learner rather than overfitting. Among fourteen prior systems surveyed, this is the only approach that unifies all seven design dimensions of affect-aware reinforcement learning-based tutoring. It combines a benchmark-validated perception layer with a deep value-based policy through a defined interface, enabling a shift from cognition-only adaptation toward tutoring systems that respond to the learner as they are, rather than solely to recorded performance.